% Source: Solutions/Srijit-PS2-Q12-15.pdf, p. 1 (typed). Restated problem statement omitted.
\submittedby{Srijit Bhadra}

\emph{Solution:} We know,
$\Int(A)\cap\Bd(A)=\Int(A)\cap\bar{A}\cap\overline{X\setminus A}$. For the sake of
contradiction, assume there is a point
$x\in\Int(A)\cap\Bd(A)\implies x\in\Int(A)\cap\bar{A}\cap\overline{X\setminus A}$. As
$x\in\overline{X\setminus A}$, every neighbourhood of $x$ must intersect $X\setminus A$.
% Corrected 2026-11-10 at the user's request: the PDF has a closure bar over X \ A here.
Thus, $\Int(A)\cap(X\setminus A)\neq\emptyset$. But this is impossible as
$\Int(A)\subseteq A$ ($\Rightarrow\!\Leftarrow$). Thus, $\Int(A)\cap\Bd(A)=\emptyset$.

For the second part, let $x\in\bar{A}$. Again, this happens if and only if every neighborhood of
$x$ intersects $A$. Now, either there exists a neighborhood $U$ of $x$ such that
$x\in U\subseteq A$, or there is no such neighborhood. In the first case,
$x\in U\subseteq A\implies x\in\Int(A)$ since $\Int(A)$ is the largest open set contained in
$A$. In the second case, every neighborhood $V$ of $x$ intersects
$X\setminus A\implies x\in\overline{X\setminus A}\implies x\in\Bd(A)$. Combining the two cases,
$x\in\Int(A)\cup\Bd(A)$. Since $x\in\bar{A}$ was chosen arbitrarily,
$\bar{A}\subseteq\Int(A)\cup\Bd(A)$.

For the other way inclusion, notice that $\Int(A)\subseteq A\subseteq\bar{A}$. Also,
$\Bd(A)=\bar{A}\cap\overline{X\setminus A}\subseteq\bar{A}$. Therefore,
$\Int(A)\cup\Bd(A)\subseteq\bar{A}$. Hence, we conclude that $\bar{A}=\Int(A)\cup\Bd(A)$.
\hfill$\square$
